% ch4_answers.tex -- Quick check answers for Module 7, Chapter 4.
\begin{chapteranswers}{4}{Light and Special Relativity}
  % ---- 4.1 Frames and postulates
  \item \final{(A)}, the cyclist on the straight, flat road. She moves in a straight line at a
        constant speed, so her velocity is constant and she is not accelerating. (B) changes
        direction, (C) slows down and (D) moves in a circle, so they are all accelerating
        (non-inertial).
  \item (1) All inertial frames of reference are equivalent. The laws of physics are the same in
        every inertial frame. No experiment can show which one is ``really'' at rest.
        (2) The speed of light in a vacuum is an absolute constant. Every inertial observer
        measures $c = \qty{3.00e8}{m.s^{-1}}$, whatever the motion of the source or the observer.
  \item Postulate 2 is about light in a \emph{vacuum}. In water, light interacts with the water
        molecules and travels more slowly: $n_\mathrm{x} = \dfrac{c}{v_\mathrm{x}} = \dfrac{\num{3.00e8}}{\num{2.25e8}}
        = 1.33$. Every inertial observer still measures \qty{3.00e8}{m.s^{-1}} for light in a
        vacuum, so the slower speed in water does not break the postulate.
  % ---- 4.2 Time dilation and length contraction
  \item The crew's clock is at every flash, so $t_0 = \qty{30.0}{s}$, and
        $\dfrac{v^2}{c^2} = 0.600^2 = 0.360$. Then
        $t = \dfrac{\qty{30.0}{s}}{\sqrt{1 - 0.360}} = \dfrac{\qty{30.0}{s}}{0.800}
        = \qty{37.5}{s}$. Answer: \final{$t = \qty{37.5}{s}$} (longer than $t_0$, as expected).
  \item The strike at end~A happens first. If the train moves left, end~A is now the front.
        The train carries Peter towards the flash from A, so it reaches him first. In his own
        frame he is at rest halfway between the ends, and both flashes travel at $c$ over equal
        distances. So, for Peter, the strike at A really happened first.
  \item (a) The station sees the shuttle move, so it measures a contracted length:
        $l = \qty{120}{m} \times \sqrt{1 - 0.60^2} = \qty{120}{m} \times 0.80 = \qty{96}{m}$.
        Answer: \final{\qty{96}{m}}. (b) The pilot is at rest relative to the shuttle, so she
        measures the proper length, \final{\qty{120}{m}}.
  \item (a) The Earth--star distance is a proper length (both are at rest in the Earth frame):
        $l = \qty{4.2}{\lightyear} \times \sqrt{1 - 0.90^2} = \qty{4.2}{\lightyear} \times 0.436
        = \qty{1.83}{\lightyear}$. Answer: \final{\qty{1.8}{\lightyear}}.
        (b) For the probe, the star comes towards it at $0.90c$ over \qty{1.83}{\lightyear}:
        $t_0 = \dfrac{1.83}{0.90} = \qty{2.03}{\yr}$. Answer: \final{\qty{2.0}{\yr}}.
        Check with time dilation: Earth time $= \dfrac{4.2}{0.90} = \qty{4.67}{\yr}$, and
        $4.67 \times 0.436 = \qty{2.04}{\yr} \approx \qty{2.0}{\yr}$ \cmark\ (the small
        difference from 2.03 comes from rounding).
  % ---- 4.3 Relativistic momentum
  \item (a) $v = 0.80 \times \qty{3.00e8}{m.s^{-1}} = \qty{2.40e8}{m.s^{-1}}$ and
        $\sqrt{1 - 0.80^2} = 0.60$, so
        $p_\mathrm{v} = \dfrac{\qty{1.5e5}{kg} \times \qty{2.40e8}{m.s^{-1}}}{0.60}
        = \qty{6.0e13}{kg.m.s^{-1}}$. Answer: \final{\qty{6.0e13}{kg.m.s^{-1}}}.
        (b) The camera moves with the probe, so for it $v = 0$ and the probe's momentum is
        \final{zero}.
  \item $p_\mathrm{v} = 2m_0v$ means $\dfrac{1}{\sqrt{1 - v^2/c^2}} = 2$. Flip both sides
        upside down: $\sqrt{1 - v^2/c^2} = \dfrac{1}{2} = 0.5$. Square both sides: $1 - \dfrac{v^2}{c^2} = 0.25$, so
        $\dfrac{v^2}{c^2} = 0.75$ and $\dfrac{v}{c} = \sqrt{0.75} = 0.866$. Answer:
        \final{$v = 0.866c = \qty{2.60e8}{m.s^{-1}}$}.
  \item A magnet bends a charged particle into a circle of radius $r = \dfrac{p}{qB}$, so for a
        fixed ring the field needed is $B = \dfrac{p}{qr}$. Near $c$ the speed hardly changes,
        but $p_\mathrm{v} = \dfrac{m_0v}{\sqrt{\left(1 - v^2/c^2\right)}}$ keeps growing quickly
        as energy is added. So $B$ must be raised in step with $p_\mathrm{v}$. If it were not, the
        radius would grow and the particles would hit the walls of the ring.
  \item Rearrange $r = \dfrac{p_\mathrm{v}}{qB}$: multiply both sides by $B$, then divide both
        sides by $r$: $B = \dfrac{p_\mathrm{v}}{qr}$. Here
        $qr = \qty{1.602e-19}{C} \times \qty{25.0}{m} = \qty{4.005e-18}{C.m}$.\par
        (a) $\sqrt{1 - 0.990^2} = \sqrt{0.0199} = 0.1411$ and
        $v = \qty{2.97e8}{m.s^{-1}}$, so
        $p_\mathrm{v} = \dfrac{\qty{1.673e-27}{kg} \times \qty{2.97e8}{m.s^{-1}}}{0.1411}
        = \qty{3.522e-18}{kg.m.s^{-1}}$ and
        $B = \dfrac{\num{3.522e-18}}{\num{4.005e-18}}\,\si{T}$. Answer:
        \final{$B = \qty{0.879}{T}$}.\par
        (b) $\sqrt{1 - 0.999^2} = \sqrt{0.001999} = 0.04471$ and
        $v = \qty{2.997e8}{m.s^{-1}}$, so
        $p_\mathrm{v} = \dfrac{\qty{1.673e-27}{kg} \times \qty{2.997e8}{m.s^{-1}}}{0.04471}
        = \qty{1.121e-17}{kg.m.s^{-1}}$ and
        $B = \dfrac{\num{1.121e-17}}{\num{4.005e-18}}\,\si{T}$. Answer:
        \final{$B = \qty{2.80}{T}$}.\par
        The speed rises by less than 1\,\% ($0.999 \div 0.990 = 1.009$), but the field must
        rise about 3.2 times ($2.800 \div 0.8795 = 3.18$).
  % ---- 4.4 E = mc^2
  \item (a) $P = \dfrac{\Delta E}{\Delta t}$ and $\Delta E = \Delta m\,c^2$, so
        $\dfrac{\Delta m}{\Delta t} = \dfrac{P}{c^2}
        = \dfrac{\qty{3.0e9}{W}}{(\qty{3.00e8}{m.s^{-1}})^2} = \qty{3.3e-8}{kg.s^{-1}}$.
        (b) In one year: $\qty{3.33e-8}{kg.s^{-1}} \times \qty{3.16e7}{s} = \qty{1.05}{kg}$.
        Answers: \final{\qty{3.3e-8}{kg}} each second and \final{about \qty{1.1}{kg}} each year.
  \item All of the rest mass becomes energy, and the two photons share it equally. So each
        photon carries the rest energy of one proton:
        $E = m_\mathrm{p}c^2 = \qty{1.673e-27}{kg} \times (\qty{3.00e8}{m.s^{-1}})^2
        = \qty{1.506e-10}{J}$. In MeV: $\dfrac{\qty{1.506e-10}{J}}{\qty{1.602e-13}{J}}
        = \qty{940}{\mega\eV}$. Answer: \final{\qty{1.51e-10}{J} = \qty{940}{\mega\eV}} each.
        (The accurate value is \qty{938}{\mega\eV}; the data-sheet $c$ is rounded.)
  \item $\Delta m = \dfrac{\Delta E}{c^2} = \dfrac{\qty{3.0e7}{J}}{\qty{9.00e16}{m^2.s^{-2}}}
        = \qty{3.3e-10}{kg}$. Answer: \final{\qty{3.3e-10}{kg}}. This is only about 3 parts in
        10 billion of the coal's mass. Even if the coal and its oxygen were sealed in a container,
        no balance could detect such a tiny fraction of the total mass. It is known only by
        calculation.
  \item $\Delta m = (2.014102 + 3.016049) - (4.002603 + 1.008665) = 5.030151 - 5.011268
        = \qty{0.018883}{\amu}$. Then $E = 0.018883 \times \qty{931.5}{\mega\eV}
        = \qty{17.59}{\mega\eV}$, and $17.59 \times \qty{1.602e-13}{J} = \qty{2.818e-12}{J}$.
        Answer: \final{\qty{17.6}{\mega\eV} = \qty{2.82e-12}{J}}.
  % ---- 4.5 Key experiments
  \item Muons of one narrow speed range (near $0.995c$) were counted at two heights. On a
        mountain \qty{1.9}{km} up there were about 563 per hour; at sea level there were 408 per
        hour. The trip down
        takes about \qty{6.4}{\micro\second} in the Earth frame: nearly 3 lifetimes. Without
        relativity only about 30 per hour should survive. Far more survive, because the muons'
        clocks run slow (a time-dilation factor of about 9). So the counts are evidence for
        time dilation.
  \item Use the Earth's centre as an (almost) inertial frame. The ground moves east as the Earth
        spins. The eastward plane's speed adds to this spin, so it moves fastest of all. The
        faster a clock moves, the slower it runs (time dilation). So the eastward clocks ran
        slower than the ground clocks and lost time. (Height made them gain a little, but the
        speed effect was bigger: $-184 + 144 = \qty{-40}{ns}$ predicted.)
  \item (a) $\dfrac{v}{c} = \dfrac{250}{\num{3.00e8}} = \num{8.33e-7}$, so
        $\dfrac{v^2}{c^2} = (\num{8.33e-7})^2 = \num{6.94e-13}$. Answer: \final{\num{6.94e-13}}.
        (b) $t = 48.0 \times \qty{3600}{s} = \qty{1.728e5}{s}$, and
        $\dfrac{v^2}{2c^2} = \num{6.94e-13} \div 2 = \num{3.47e-13}$. Shortcut:
        $t - t_0 \approx t \times \dfrac{v^2}{2c^2} = \qty{1.728e5}{s} \times \num{3.47e-13}
        = \qty{6.0e-8}{s}$. Answer: \final{about \qty{60}{ns}}. This is the same size as the
        $-184$ and $+96$ ns in the Hafele--Keating table.
  \item The distant fading is stretched by $1 + z = 1 + 0.50 = 1.50$. So the nearby time is
        $\dfrac{30}{1.50} = 20$ days. Answer: \final{20 days}.
\end{chapteranswers}
